\SourceLeaf{415}
\[
\begin{aligned}\clemn\varphi&=1-\varphi\varphi+\dfrac12\varphi^4-\dfrac3{10}\varphi^6\\[6pt]&\quad{}+\dfrac7{40}\varphi^8-\dfrac{61}{600}\varphi^{10}+\dfrac{71}{1200}\varphi^{12}\mathord{.}\mathord{.}\mathord{.}\mathord{.}\end{aligned}
\]
\[
\begin{aligned}
\left.\begin{matrix}p\varphi\\q\varphi\end{matrix}\right\}
&=1\mp\dfrac12\varphi\varphi-\dfrac1{24}\varphi^4\mp\dfrac1{240}\varphi^6\\[6pt]&\quad{}+\dfrac{17}{40320}\varphi^8\mp\dfrac1{403200}\varphi^{10}+\dfrac{37}{159667200}\varphi^{12}\\[6pt]
&\quad\pm\dfrac{113}{4151347200}\varphi^{14}+\dfrac{4171}{6974263296000}\varphi^{16}\mathord{.}\mathord{.}
\end{aligned}
\]
Formulae pro $P$, $Q$, $p$, $q$ in infinitum continuatae quavis convergentia data citius convergunt, formula autem pro $\slemn\varphi$ diverget, si $\varphi$ ponetur $>\frac{\varpi}{\surd2}$ sive $\varphi^4>\frac{\varpi^4}{4}$, formula autem pro $\clemn\varphi$ diverget, si $\varphi>\varpi$
\par\begin{center}\rule{.22\linewidth}{.3pt}\end{center}
\par\Needspace{8\baselineskip}\pdfbookmark[0]{Einige neue Formeln die Lemniscatischen Functionen betreffend}{lemn-new-formulas}\begin{center}\textit{Einige neue Formeln die Lemniscatischen Functionen betreffend.}\end{center}\selectlanguage{german}
Es sei $\displaystyle\int\frac{\mathrm{d}x}{\surd(1-x^4)}=\varphi$, oder
\[
\begin{aligned}\varphi&=x+\dfrac12.\dfrac15x^5+\dfrac{1.3}{2.4}.\dfrac19x^9\\[6pt]&\quad{}+\dfrac{1.3.5}{2.4.6}\dfrac1{13}x^{13}\mathord{.}\mathord{.}\end{aligned}
\]

Man hat dann
\[
\begin{aligned}\varphi\varphi&=xx+\dfrac35.\dfrac13x^6+\dfrac{3.7}{5.9}.\dfrac15x^{10}\\[6pt]&\quad{}+\dfrac{3.7.11}{5.9.13}.\dfrac17x^{14}.\end{aligned}
\]
Es sei $x=\slemn\varphi$, $y=\slemn\psi$, $z=\slemn(\varphi+\psi)$

so hat man
\[
\begin{aligned}z&=\dfrac{x\surd(1-y^4)+y\surd(1-x^4)}{1+xxyy}\\[6pt]&=\dfrac{xx-yy}{x\surd(1-y^4)-y\surd(1-x^4)}\end{aligned}
\]
\[
\begin{aligned}\surd\dfrac{1-zz}{1+zz}&=\dfrac{-2xy+\surd(1-x^4)\surd(1-y^4)}{1+xx+yy-xxyy}\\[6pt]&=\dfrac{1-xx-yy-xxyy}{2xy+\surd(1-x^4)\surd(1-y^4)}\end{aligned}
\]
\[
\begin{aligned}\surd\dfrac{1-z^4}{zz}&=\dfrac{\begin{gathered}x(1+y^4)\surd(1-x^4)\\[3pt]{}-y(1+x^4)\surd(1-y^4)\end{gathered}}{(1+xxyy)(xx-yy)}\end{aligned}
\]
\[
\begin{aligned}\surd(1-z^4)&=\dfrac{\begin{gathered}(1-xxyy)\surd(1-x^4)\surd(1-y^4)\\[3pt]{}-2xy(xx+yy)\end{gathered}}{(1+xxyy)^2}\end{aligned}
\]
\[
\begin{aligned}\surd(1-zz)&=\dfrac{\begin{gathered}\surd(1-xx)\surd(1-yy)\\[3pt]{}-xy\surd(1+xx)\surd(1+yy)\end{gathered}}{1+xxyy}\end{aligned}
\]
\[
\begin{aligned}\surd(1+zz)&=\dfrac{\begin{gathered}\surd1+xx)\surd(1+yy)\\[3pt]{}+xy\surd(1-xx)\surd(1-yy)\end{gathered}}{1+xxyy}\end{aligned}
\]

\SourceLeaf{416}
Die einfachste Manier $\slemn\varphi$ in eine Reihe nach Potenzen von $\varphi$ zu entwickeln scheint folgende zu sein: man hat
\[
\slemn(1+i)\varphi=\dfrac{(1+i)\slemn\varphi}{\surd(1-\slemn\varphi^4)}
\]
setzt man also
\[
(\slemn\varphi)^{-2}=\varphi^{-2}(1+\alpha\varphi^4+\beta\varphi^8+\gamma\varphi^{12}\mathord{.}\mathord{.}\mathord{.}\mathord{.})
\]
und folglich
\[
\begin{aligned}(\slemn(1+i)\varphi)^{-2}&=-\tfrac12i\varphi^{-2}(1-4\alpha\varphi^4+16\beta\varphi^8\\[6pt]&\quad{}-64\gamma\varphi^{12}\mathord{.}\mathord{.}\mathord{.})\end{aligned}
\]
also
\[
\begin{aligned}
\slemn\varphi^2&=-2i(\slemn(1+i)\varphi)^{-2}+(\slemn\varphi)^{-2}\\[6pt]
&=5\alpha\varphi\varphi-15\beta\varphi^6+65\gamma\varphi^{10}-255\delta\varphi^{14}\mathord{.}\mathord{.}\mathord{.}
\end{aligned}
\]
man hat also
\[
\begin{aligned}(1+\alpha t+\beta t^2+\gamma t^3+\delta t^4+\mathord{.}\mathord{.}\mathord{.})\\[6pt]&\quad(5\alpha-15\beta t+65\gamma t^2-255\delta t^3\\[5pt]&\quad{}+1025\varepsilon t^4\mathord{.}\mathord{.}\mathord{.})=1\end{aligned}
\]
hieraus
\[
\begin{aligned}5\alpha=1\\[5pt]&\alpha=\dfrac15\\[5pt]&\quad=\dfrac15\end{aligned}
\]
\[
\begin{aligned}3\beta=\alpha\alpha\\[5pt]&\beta=\dfrac1{75}\\[5pt]&\quad=\dfrac1{3.25}\\[5pt]&\quad=\dfrac1{15}\alpha\end{aligned}
\]
\[
\begin{aligned}13\gamma=2\alpha\beta\\[5pt]&\gamma=\dfrac2{4875}\\[5pt]&\quad=\dfrac2{3.5^3.13}\\[5pt]&\quad=\dfrac2{65}\beta\end{aligned}
\]
\[
\begin{aligned}51\delta=14\alpha\gamma-3\beta\beta\\[5pt]&\delta=\dfrac1{82875}\\[5pt]&\quad=\dfrac1{3.5^3.13.17}\\[5pt]&\quad=\dfrac1{34}\gamma\end{aligned}
\]
\[
\begin{aligned}205\varepsilon=50\alpha\delta-10\beta\gamma\\[5pt]&\varepsilon=\dfrac2{6215625}\\[5pt]&\quad=\dfrac2{9.5^5.13.17}\\[5pt]&\quad=\dfrac2{75}\delta\end{aligned}
\]
\[
\begin{aligned}819\zeta=206\alpha\varepsilon-54\beta\delta+13\gamma\gamma\\[5pt]&\zeta=\dfrac2{242409375}\\[5pt]&\quad=\dfrac2{3^3.5^5.13^2.17}\\[5pt]&\quad=\dfrac1{39}\varepsilon\end{aligned}
\]
\[
\begin{aligned}3277\eta=818\alpha\zeta-202\beta\varepsilon+38\gamma\delta\\[5pt]&\eta=\dfrac4{19527421875}\\[5pt]&\quad=\dfrac4{3.5^7.13^2.17.29}\\[5pt]&\quad=\dfrac{18}{725}\zeta\end{aligned}
\]
\[
\begin{aligned}13107\theta=3278\alpha\eta-822\beta\zeta+218\gamma\varepsilon-51\delta\delta,\\[5pt]&\theta=\dfrac{223}{44815433203125}\\[5pt]&\quad=\dfrac{223}{3^4.5^8.13^2.17^2.29}\\[5pt]&\quad=\dfrac{223}{9180}\eta\end{aligned}
\]
Die Grenze des Verhältnisses zweier aufeinander folgender Glieder ist
\[
(2{,}6220\mathord{.}\mathord{.})^4:1=47{,}27:1
\]

\SourceLeaf{417}

Sehr nahe ist
\[
\zeta=\dfrac{92}{(2\varpi)^{24}},\qquad
\eta=\dfrac{108}{(2\varpi)^{28}},\qquad
\theta=\dfrac{124}{(2\varpi)^{32}}
\]
\[
\begin{aligned}
\log P&=\log\varphi-\dfrac1{12}\alpha\varphi^4-\dfrac1{56}\beta\varphi^8
\\[6pt]&\quad{}-\dfrac1{132}\gamma\varphi^{12}-\dfrac1{240}\delta\varphi^{16}
\\[6pt]&\quad{}-\dfrac1{380}\varepsilon\varphi^{20}-\dfrac1{552}\zeta\varphi^{24}\mathord{.}\mathord{.}\mathord{.}\\[6pt]
&=\log\varphi-\dfrac1{60}\varphi^4-\dfrac1{4200}\varphi^8
\\[6pt]&\quad{}-\dfrac1{321750}\varphi^{12}-\dfrac1{19890000}\varphi^{16}
-\dfrac1{1180968750}\varphi^{20}\\[6pt]
&\quad-\dfrac{22}{1756255921875}\varphi^{24}-\mathord{.}\mathord{.}
\end{aligned}
\]
\[
\begin{aligned}
\log\operatorname{sin\,lemn}\varphi
&=\log\varphi-\dfrac36\alpha\varphi^4+\dfrac7{28}\beta\varphi^8
\\[6pt]&\quad{}-\dfrac{33}{66}\gamma\varphi^{12}+\dfrac{127}{120}\delta\varphi^{16}
\\[6pt]&\quad{}-\dfrac{513}{190}\varepsilon\varphi^{20}+\dfrac{2047}{276}\zeta\varphi^{24}\mathord{.}\mathord{.}\\[6pt]
&=\log\varphi-\dfrac1{10}\varphi^4+\dfrac1{300}\varphi^8
\\[6pt]&\quad{}-\dfrac1{4875}\varphi^{12}+\dfrac{127}{9945000}\varphi^{16}
-\dfrac3{3453125}\varphi^{20}\\[6pt]
&\quad+\dfrac{89}{1454456250}\varphi^{24}-\mathord{.}\mathord{.}
\end{aligned}
\]
Die Coëfficienten $\alpha$, $\beta$, $\gamma$ u.\,s.\,f. lassen sich auch vermittelst folgender Gleichung bestimmen
\[
\begin{aligned}1+\alpha t+\beta tt+\gamma t^3+\mathord{.}\mathord{.}\mathord{.}
\\[6pt]&=\left(\begin{gathered}1+\dfrac3{2.3}\alpha t+\dfrac7{8.7}\beta tt
\\[4pt]{}+\dfrac{33}{32.11}\gamma t^3+\dfrac{127}{128.15}\delta t^4\mathord{.}\mathord{.}\mathord{.}\end{gathered}\right)^2\end{aligned}
\]
Die bequemste Art $\log\operatorname{cos\,lemn}\varphi$ in eine Reihe zu entwickeln ist folgende. Es ist
\[
\begin{aligned}\dfrac{\mathrm d\log\operatorname{cos\,lemn}\varphi}{\mathrm d\varphi}
&=-\dfrac{2\mathrm d\varphi'}{\mathrm d\log\operatorname{sin\,lemn}\varphi}
\\[7pt]&=-(1-i)\operatorname{sin\,lemn}(1+i)\varphi\end{aligned}
\]
hieraus ergibt sich
\[
\begin{aligned}
\log\operatorname{cos\,lemn}\varphi
&=-\varphi\varphi-\dfrac2{15}\varphi^6-\dfrac2{75}\varphi^{10}
\\[6pt]&\quad{}-\dfrac{44}{6825}\varphi^{14}-\dfrac{422}{248625}\varphi^{18}
\\[6pt]&\quad{}-\dfrac{6428}{13673375}\varphi^{22}\\[6pt]
&\quad-\dfrac{6044}{44890625}\varphi^{26}
-\dfrac{20824792}{527240390625}\varphi^{30}-\mathord{.}\mathord{.}
\end{aligned}
\]
\[
P\dfrac{\mathrm d^4P}{\mathrm d\varphi^4}
-4\dfrac{\mathrm dP}{\mathrm d\varphi}\, .\,\dfrac{\mathrm d^3P}{\mathrm d\varphi^3}
+3\left(\dfrac{\mathrm d\mathrm dP}{\mathrm d\varphi^2}\right)^2=2PP
\]
\[
\begin{aligned}
P\varphi&=\varphi-\dfrac2{1\mathord{.}\mathord{.}\mathord{.}\mathord{.}5}\varphi^5
-\dfrac{36}{1\mathord{.}\mathord{.}\mathord{.}\mathord{.}9}\varphi^9
\\[6pt]&\quad{}+\dfrac{552}{1\mathord{.}\mathord{.}\mathord{.}\mathord{.}13}\varphi^{13}
+\dfrac{5136}{1\mathord{.}\mathord{.}\mathord{.}\mathord{.}17}\varphi^{17}
\\[6pt]&\quad{}+\dfrac{5146848}{1\mathord{.}\mathord{.}\mathord{.}\mathord{.}21}\varphi^{21}-\mathord{.}\mathord{.}\\[6pt]
&=\varphi-\dfrac1{4.3.5}\varphi^5-\dfrac1{32.9.5.7}\varphi^9
\\[6pt]&\quad{}+\dfrac{23}{128.81.25.7.11.13}\varphi^{13}
\\[6pt]&\quad{}+\dfrac{107}{2048.243.125.49.11.13.17}\varphi^{17}\\[6pt]
&\quad+\dfrac{23.37}{8192.729.625.49.11.13.17.19}\varphi^{21}-\mathord{.}\mathord{.}
\end{aligned}
\]
\SourceLeaf{418}
\[
\begin{aligned}
P\psi\varpi&=+2.6220575\,\psi\\
&\phantom{{}={}}-2.0656648\,\psi^5\\
&\phantom{{}={}}-0.5811918\,\psi^9\\
&\phantom{{}={}}+0.0245475\,\psi^{13}\\
&\phantom{{}={}}+0.0001890\,\psi^{17}\\
&\phantom{{}={}}+0.0000620\,\psi^{21}
\end{aligned}
\]
\[
P\tfrac12\varpi=+1.2453446
={}^{4}\!\surd\tfrac12\, .\,e^{\dfrac18\pi}
\]
\[
\dfrac{\mathrm d\mathrm d\log Q}{\mathrm d\varphi^2}=\dfrac{PP}{QQ}
\]
Die Halbirung geschieht sehr bequem so
\[
\begin{aligned}
\operatorname{cos\,lemn}\varphi&=\operatorname{tang}u\\[8pt]
\operatorname{sin\,lemn}\varphi&=\surd\cos2u\\[8pt]
\operatorname{cos\,lemn}\varphi&=\surd\operatorname{tang}\tfrac12(45^0+u)\\[8pt]
\operatorname{sin\,lemn}\varphi&=\surd\operatorname{tang}\tfrac12(45^0-u)
\end{aligned}
\]
\[
\begin{aligned}
\operatorname{sin\,lemn}(\varphi+i\psi)&=r(\cos v+i\sin v)\\[8pt]
\operatorname{sin\,lemn}\varphi^2&=\operatorname{tang}\Phi\\[8pt]
\operatorname{sin\,lemn}\psi^2&=\operatorname{tang}\Psi\\[8pt]
r&=\surd\operatorname{tang}(\Phi+\Psi)\\[8pt]
\operatorname{tang}v&=\surd\dfrac{\operatorname{tang}2\Psi}{\operatorname{tang}2\Phi}
\end{aligned}
\]
\[
\begin{aligned}
\operatorname{sin\,lemn}(\varphi+\psi)
&=\dfrac{\surd(\cos2\Phi\sin2\Psi)+\surd(\sin2\Phi\cos2\Psi)}{\cos(\Phi-\Psi)\, .\,\surd2}\\[8pt]
&=\dfrac{\sin(\Psi-\Phi)\, .\,\surd2}{\surd(\cos2\Phi\sin2\Psi)-\surd(\sin2\Phi\cos2\Psi)}
\end{aligned}
\]
Es sei
\[
\operatorname{sin\,lemn}A=x,\qquad
\operatorname{sin\,lemn}B=y,\qquad
\operatorname{sin\,lemn}\dfrac{A+B}{1+i}=z
\]
so ist
\[
\begin{aligned}(1+i)z&=\dfrac{\begin{gathered}\surd(1+xx)(1-yy)\\[3pt]{}-\surd(1-xx)(1+yy)\end{gathered}}{x-y}\end{aligned}
\]
\[
\begin{aligned}\dfrac{1-i}{z}&=\dfrac{\begin{gathered}\surd1+xx)(1-yy)\\[3pt]{}+\surd(1-xx)(1+yy)\end{gathered}}{x+y}\end{aligned}
\]
